% GENERATED FILE. Edit canonical lessons, then run npm run generate:books.
% canonical-source-hash: fnv1a64:a355d7250821bf8c
% canonical-lessons: TE-C148-kapadu, TE-C148-kalcu, TE-C148-arpu, TE-C148-padu, TE-C148-ammu

\chapter{To protect, To burn, To put out, To sing, To sell}
\label{ch:te-to-protect-to-burn-to-put-out-to-sing-to-sell}
\hlchaptermodality{148}

\begin{chapteropening}
\textbf{By the end of this chapter:} \emph{I can say to protect, to burn, to put out, to sing and to sell.}

Complete the last lesson of chapter 148: i can say to protect, to burn, to put out, to sing and to sell.
\end{chapteropening}

\section[kāpāḍu]{\te{కాపాడు} (kāpāḍu) — to protect, to save}

\label{lesson:TE-C148-kapadu}



\begin{quote}
Before the new one: say the Telugu for to win, then the Telugu for to hide.
\end{quote}

\subsection*{You'll want to know: \te{కాపాడు}}
\textbf{\te{కాపాడు}} — \emph{kāpāḍu} — \textquotedblleft{}to protect, to save\textquotedblright{}.

\begin{grammarlens}[title={What you've built}]
One more word for this chapter.
\end{grammarlens}

\subsection*{Your turn}
\begin{itemize}
\raggedright
  \item \emph{Say it:} \emph{kāpāḍu}
  \item \emph{Say it:} \emph{kāpāḍu}, once more
  \item \emph{Say it:} say \emph{kāpāḍu}
  \item \emph{Recall:} say the Telugu for to win, then the Telugu for to hide, then say \emph{kāpāḍu} again
\end{itemize}

\begin{tcolorbox}[breakable,colback=teal!4,colframe=teal!35!black,title={Before you move on}]
What does \te{కాపాడు} mean? (\textquotedblleft{}To protect, to save\textquotedblright{}.) Say it once more.
\end{tcolorbox}
\section[kālcu]{\te{కాల్చు} (kālcu) — to burn}

\label{lesson:TE-C148-kalcu}



\begin{quote}
Before the new one: say the Telugu for to hide, then the Telugu for to protect.
\end{quote}

\subsection*{You'll want to know: \te{కాల్చు}}
\textbf{\te{కాల్చు}} — \emph{kālcu} — \textquotedblleft{}to burn\textquotedblright{}.

\begin{grammarlens}[title={What you've built}]
One more word for this chapter.
\end{grammarlens}

\subsection*{Your turn}
\begin{itemize}
\raggedright
  \item \emph{Say it:} \emph{kālcu}
  \item \emph{Say it:} \emph{kālcu}, once more
  \item \emph{Say it:} say \emph{kālcu}
  \item \emph{Recall:} say the Telugu for to hide, then the Telugu for to protect, then say \emph{kālcu} again
\end{itemize}

\begin{tcolorbox}[breakable,colback=teal!4,colframe=teal!35!black,title={Before you move on}]
What does \te{కాల్చు} mean? (\textquotedblleft{}To burn\textquotedblright{}.) Say it once more.
\end{tcolorbox}
\section[ārpu]{\te{ఆర్పు} (ārpu) — to put out}

\label{lesson:TE-C148-arpu}



\begin{quote}
Before the new one: say the Telugu for to protect, then the Telugu for to burn.
\end{quote}

\subsection*{You'll want to know: \te{ఆర్పు}}
\textbf{\te{ఆర్పు}} — \emph{ārpu} — \textquotedblleft{}to put out\textquotedblright{}.

\begin{grammarlens}[title={What you've built}]
One more word for this chapter.
\end{grammarlens}

\subsection*{Your turn}
\begin{itemize}
\raggedright
  \item \emph{Say it:} \emph{ārpu}
  \item \emph{Say it:} \emph{ārpu}, once more
  \item \emph{Say it:} say \emph{ārpu}
  \item \emph{Recall:} say the Telugu for to protect, then the Telugu for to burn, then say \emph{ārpu} again
\end{itemize}

\begin{tcolorbox}[breakable,colback=teal!4,colframe=teal!35!black,title={Before you move on}]
What does \te{ఆర్పు} mean? (\textquotedblleft{}To put out\textquotedblright{}.) Say it once more.
\end{tcolorbox}
\section[pāḍu]{\te{పాడు} (pāḍu) — to sing}

\label{lesson:TE-C148-padu}



\begin{quote}
Before the new one: say the Telugu for to burn, then the Telugu for to put out.
\end{quote}

\subsection*{You'll want to know: \te{పాడు}}
\textbf{\te{పాడు}} — \emph{pāḍu} — \textquotedblleft{}to sing\textquotedblright{}.

\begin{grammarlens}[title={What you've built}]
One more word for this chapter.
\end{grammarlens}

\subsection*{Your turn}
\begin{itemize}
\raggedright
  \item \emph{Say it:} \emph{pāḍu}
  \item \emph{Say it:} \emph{pāḍu}, once more
  \item \emph{Say it:} say \emph{pāḍu}
  \item \emph{Recall:} say the Telugu for to burn, then the Telugu for to put out, then say \emph{pāḍu} again
\end{itemize}

\begin{tcolorbox}[breakable,colback=teal!4,colframe=teal!35!black,title={Before you move on}]
What does \te{పాడు} mean? (\textquotedblleft{}To sing\textquotedblright{}.) Say it once more.
\end{tcolorbox}
\section[ammu]{\te{అమ్ము} (ammu) — to sell}

\label{lesson:TE-C148-ammu}



\begin{quote}
Before the new one: say the Telugu for to put out, then the Telugu for to sing.
\end{quote}

\subsection*{You'll want to know: \te{అమ్ము}}
\textbf{\te{అమ్ము}} — \emph{ammu} — \textquotedblleft{}to sell\textquotedblright{}.

\begin{grammarlens}[title={What you've built}]
One more word for this chapter.
\end{grammarlens}

\subsection*{Your turn}
\begin{itemize}
\raggedright
  \item \emph{Say it:} \emph{ammu}
  \item \emph{Say it:} \emph{ammu}, once more
  \item \emph{Say it:} say \emph{ammu}
  \item \emph{Recall:} say the Telugu for to put out, then the Telugu for to sing, then say \emph{ammu} again
\end{itemize}

\begin{tcolorbox}[breakable,colback=teal!4,colframe=teal!35!black,title={Before you move on}]
What does \te{అమ్ము} mean? (\textquotedblleft{}To sell\textquotedblright{}.) Say it once more.
\end{tcolorbox}
